\chapter{Исследовательский раздел}

В данном разделе описано исследование качества результатов работы
модифицированного метода ведения диалога роботом Ф-2 с учётом персональных
черт робота.

\section{Постановка и условия исследования}

Согласно формуле~(\ref{eq:personality_profile}), роботу назначается профиль, который задаётся термами персональных черт.
Дано 30 граней --- по 6 граней на каждую черту.
Значение каждой грани описывается одним из трёх возможных термов: низкое, среднее, высокое значение.

Так, базовый метод устроен так, чтобы сформировать условно положительного собеседника без крайних проявлений черт характера.
В терминах предложенного модифицированного метода такой базовый профиль робота можно представить через 30 граней, имеющих значение ``среднее''.

Интерес представляет установление посредством опроса того, различают ли респонденты в конкретных реакциях степени проявления различных граней.

\section{Описание анкеты и респондентов}

Опрос проводился с помощью электронной анкеты.
Во вводной части анкеты респондентам объяснялась задача исследования: определить, какие реакции робота в заданных диалоговых ситуациях лучше соответствуют предполагаемому профилю робота и конкретной степени проявления грани личности.
Также во вводной части был приведён пример заполнения: для одной ситуации показывались
несколько возможных реакций робота, после чего отдельно рассматривался выбор
реакций для низкой, средней и высокой степени проявления грани.

Основная часть анкеты состояла из 30 вопросов для каждой из 30 граней Большой пятёрки черт.
Каждый вопрос имел одинаковую структуру:

\begin{enumerate}
	\item краткое текстовое пояснение смысла рассматриваемой грани;
	\item описание одной диалоговой ситуации в виде реплики собеседника,
	адресованной роботу;
	\item набор возможных текстовых реакций робота на данную реплику;
	\item возможность выбора реакций при низкой, средней и высокой степени проявления соответствующей грани в профиле робота.
\end{enumerate}

В каждом вопросе респонденту предлагалось выбрать одну или две реакции, которые, по его мнению, лучше всего отражают поведение робота при указанной степени проявления грани.
Например, для грани <<дружелюбие>> рассматривалась ситуация, в которой собеседник сообщает роботу, что собирается помочь своему другу; респондент должен был выбрать реакции робота отдельно для низкой, средней и высокой степени дружелюбия.
Аналогичным образом были сформированы вопросы для остальных граней.

Информация о респондентах:

\begin{itemize}
	\item количество --- 19;
	\item возраст --- от 22 до 37 лет.
\end{itemize}

\section{Критерии качества}

Для оценки результатов исследования выбраны приведенные ниже критерии качества.


% ввести обозначения
%\begin{enumerate}
	%\item
\textbf{Согласованность ответов респондентов} показывает, насколько устойчиво люди соотносят предложенные реакции робота с заданной степенью проявления грани личности.
	Для каждой пары <<ситуация + степень>> согласованность определялась как доля респондентов, выбравших наиболее частый полный набор реакций.
	Дополнительно для каждой грани вычислялась средняя согласованность по трём степеням: низкой, средней и высокой.
	Расчёт выполнялся по трём критериям, согласно следующим формулам:
	\begin{equation}
		\label{eq:research_consensus}
		C = \frac{n_{\max}}{N},
	\end{equation}
	где
	\begin{itemize}
		\item \(N\) --- общее число респондентов;
		\item \(n_{\max}\) --- число респондентов, выбравших самый частый полный набор реакций для рассматриваемой пары <<ситуация + степень>>;
		\item \(C\) --- согласованность ответов респондентов для рассматриваемой пары <<ситуация + степень>>,
	\end{itemize}

	\begin{equation}
		\label{eq:research_facet_consensus}
		C_{\text{грани}} = \frac{C_{\text{низк}} + C_{\text{сред}} + C_{\text{выс}}}{3},
	\end{equation}
	где
	\begin{itemize}
		\item \(C_{\text{низк}}\) --- согласованность для низкой степени проявления грани;
		\item \(C_{\text{сред}}\) --- согласованность для средней степени проявления грани;
		\item \(C_{\text{выс}}\) --- согласованность для высокой степени проявления грани;
		\item \(C_{\text{грани}}\) --- средняя согласованность по трём степеням проявления рассматриваемой грани,
	\end{itemize}

	\begin{equation}
		\label{eq:research_level_consensus}
		C_l = \frac{1}{K_l} \sum_{j = 1}^{K_l} C_{j,l}, \quad l \in \{\text{низк}, \text{сред}, \text{выс}\},
	\end{equation}
	где
	\begin{itemize}
		\item \(l\) --- рассматриваемая степень проявления грани;
		\item \(K_l\) --- количество пар <<ситуация + степень>>, относящихся к степени \(l\);
		\item \(C_{j,l}\) --- согласованность для \(j\)-й пары <<ситуация + степень>> при степени проявления \(l\);
		\item \(C_l\) --- средняя согласованность по всем парам «ситуация + степень», относящимся к степени проявления \(l\).
	\end{itemize}

\section{Результаты исследования согласованности мнений респондентов}

На рисунках~\ref{img:research_agreement_heatmap} и~\ref{img:research_agreement_by_facet} представлены тепловые карты согласованности ответов по граням и по граням и степеням проявления.

\includesvgimage
{research_agreement_heatmap}
{f}
{H}
{0.95\textwidth}
{Согласованность ответов по граням и степеням проявления}

\includesvgimage
{research_agreement_by_facet}
{f}
{H}
{0.95\textwidth}
{Средняя согласованность ответов по граням}

Средняя согласованность по всем парам <<ситуация + степень>> составила
50~\%.
Полное совпадение ответов всех респондентов не наблюдалось ни в одном
из 90 случаев.
В 12 случаях из 90 доля согласованности была не ниже 75~\%.
В 24 случаях доля согласованности была не ниже 67~\%.

По степеням проявления получены следующие средние значения согласованности:
для низкой  --- 57~\%, для средней --- 41~\%, для высокой ---
53~\%.

Таким образом, крайние степени в среднем распознавались лучше, чем средние.
Средняя степень оказалась наиболее неоднозначной для респондентов, поскольку она слабее отделена от крайних проявлений грани.

Также было проанализировано число ситуаций в данных от числа респондентов, давших идентичный ответ.
Полученное распределение приведено на рисунке~\ref{img:research_judgment_intersections_histogram}.

\includesvgimage
{research_judgment_intersections_histogram}
{f}
{H}
{0.95\textwidth}
{Распределение частоты выбора отдельных реакций}

При построении данного распределения учитывалось не полное совпадение
набора выбранных реакций, а совпадение отдельных реакций внутри каждой
пары <<ситуация + степень>>.
Всего было выделено 389 отдельных сочетаний <<вопрос + реакция>>.
Из них 127 сочетаний, то есть 33~\%, были выбраны не менее чем 10
респондентами из 19.
Это показывает, что для значительной части реакций наблюдается
согласованное большинство.
14~\% ответов были уникальные.

В то же время 86 сочетаний, то есть 22~\%, были выбраны только одним
или двумя респондентами.
Такие реакции можно рассматривать как индивидуальные или слабо
поддержанные варианты интерпретации ситуации.
Наибольшая часть распределения приходится на промежуточную область:
155 сочетаний были выбраны от 4 до 9 респондентов.
Следовательно, данные содержат не только явно согласованные и явно
единичные реакции, но и широкий слой частично разделяемых вариантов.
Это соответствует общей картине исследования: респонденты часто
выделяли одни и те же значимые реакции, однако полное совпадение
набора ответов оставалось редким из-за возможности выбрать один или
два варианта реакции.

Грани с наибольшей согласованностью мнений об обусловленных ими реакциях приведены в таблице~\ref{tbl:best_facets}.

\begin{table}[H]
	\centering
	\small
	\caption{Грани с наибольшей согласованностью мнений об обусловленных ими реакциях}
	\label{tbl:best_facets}
	\begin{tabularx}{\textwidth}{|c|X|p{0.24\textwidth}|p{0.18\textwidth}|}
		\hline
		\textbf{N} & \textbf{Грань} & \centering\textbf{Средняя согласованность} & \centering\arraybackslash\textbf{Полное совпадение} \\
		\hline
		29 & интеллектуальная любознательность & 81~\% & 0/3 \\
		\hline
		9 & альтруизм & 79~\% & 0/3 \\
		\hline
		17 & самодисциплина & 79~\% & 0/3 \\
		\hline
		1 & дружелюбие & 70~\% & 0/3 \\
		\hline
		15 & исполнительность & 67~\% & 0/3 \\
		\hline
		7 & доверчивость & 63~\% & 0/3 \\
		\hline
		18 & осторожность & 63~\% & 0/3 \\
		\hline
		25 & воображение & 63~\% & 0/3 \\
		\hline
	\end{tabularx}
\end{table}

Наиболее согласованной оказалась грань интеллектуальной любознательности.
Также сравнительно высокая согласованность получена для альтруизма и самодисциплины.
Это показывает, что для данных граней выбранные реакции лучше различают степени проявления, чем для остальных граней.

Грани с наименьшей согласованностью мнений об обусловленных ими реакция приведены в таблице~\ref{tbl:weak_facets}.

\begin{table}[H]
	\centering
	\small
	\caption{Грани с наименьшей согласованностью мнений об обусловленных ими реакциях}
	\label{tbl:weak_facets}
	\begin{tabularx}{\textwidth}{|c|X|p{0.24\textwidth}|p{0.18\textwidth}|}
		\hline
		\textbf{N} & \textbf{Грань} & \centering\textbf{Средняя согласованность} & \centering\arraybackslash\textbf{Полное совпадение} \\
		\hline
		8 & моральность & 28~\% & 0/3 \\
		\hline
		24 & уязвимость & 28~\% & 0/3 \\
		\hline
		21 & подавленность & 30~\% & 0/3 \\
		\hline
		13 & самоэффективность & 33~\% & 0/3 \\
		\hline
		5 & поиск острых ощущений & 35~\% & 0/3 \\
		\hline
		22 & чувствительность к оценке & 35~\% & 0/3 \\
		\hline
		3 & настойчивость & 37~\% & 0/3 \\
		\hline
		6 & жизнерадостность & 37~\% & 0/3 \\
		\hline
	\end{tabularx}
\end{table}

Низкая согласованность для перечисленных граней означает, что респонденты
по-разному соотносили предложенные реакции со степенями проявления черты.
Причинами могут быть неоднозначность диалоговой ситуации, пересечение смыслов
нескольких реакций, недостаточная различимость средней и крайней степени, а также ограниченный объём выборки.
Для таких граней требуется уточнение ситуаций и формулировок реакций перед использованием в устойчивой конфигурации правил.

\section{Отбор сценариев коммуникативного поведения с привлечением респондентов}

Архитектура робота Ф-2 служит платформой для апробации различных сценариев коммуникативного поведения.
Для апробации таких сценариев респондентам предлагается оценить их в конкретных предлагаемых условиях (ситуациях) по различным критериям~\cite{volkova2024crowdsourcing}.

В данной работе использовалась следующая исследовательская процедура.

\begin{itemize}
	\item Этап 1. Для каждой грани описать правила, которые при заданных значениях степени проявления грани БПЧ модифицируют веса активированных сценариев, при этом правило может затрагивать более одного сценария из пула д-сценариев робота Ф-2.
	\item Этап 2. Составление набора ситуаций, в которых будут рассматриваться реакции в коммуникации.
	\item Этап 3. Составить анкету и предложить респондентам указать 1 или несколько сценариев, которые более всего соответствуют по их мнению реакции предполагаемой модификации робота с заданной степенью проявления одной заданной грани (в остальном значения граней считаются средними).
	\item Этап 4. Агрегировать полученные ответы, оценить согласованность их мнений. Выполнить сравнительный анализ полученных мнений, в том числе на предмет применимости первичного набора правил.
	\item Этап 5. Выполнить отбор реакций на следующих основаниях:
	\begin{itemize}
		\item единство мнений свидетельствует в пользу общего понимания респондентами различий в поведении, обусловленных значениями конкретных граней БПЧ;
		\item сценарии реакции с большим числом голосов в заданных условиях следует отбирать для дальнейшего использования при формировании коммуникативных реакций роботом-собеседником, тогда как сценарии с единичными голосами следует отбраковывать или резко понижать их вес среди активированных сценариев.
	\end{itemize}
	\item Этап 6. Сформулировать рекомендации о модификации набора правил на основании интерпретации мнений респондентов.
\end{itemize}

	%\item
	%Соответствие разработанному методу показывает, насколько часто респонденты выбирали реакции, заложенные в разработанном методе.

\section{Результаты отбора}

В рамках этапа 1 составлен набор правил.

В рамках этапа 2 составлен набор из 30 ситуаций, в которых предполагаются вариативные реакции в зависимости от значений степени проявления конкретных граней. Анкета доступна по ссылке~\cite{questionarree}.

На этапе 3 данный набор из 30 ситуаций был предложен респондентам.
	Пример вопроса представлен на рисунке~\ref{img:example_of_questionarre_answer}.

	\includeimage
	{example_of_questionarre_answer}
	{f}
	{H}
	{0.95\textwidth}
	{Согласованность ответов по граням и степеням проявления}

В рамках этапа 4 результаты были агрегированы по следующим формулам:
		\begin{equation}
			\label{eq:research_method_match}
			M = \frac{n_{\text{метод}}}{N},
		\end{equation}
		где
		\begin{itemize}
			\item \(n_{\text{метод}}\) --- число респондентов, выбравших вариант реакций, совпадающий с результатом разработанного метода;
			\item \(N\) --- общее число респондентов;
			\item \(M\) --- совпадение с результатом разработанного метода для рассматриваемой пары <<ситуация + степень>>.
		\end{itemize}

		\begin{equation}
			\label{eq:research_level_method_match}
			M_l = \frac{1}{K_l} \sum_{j = 1}^{K_l} M_{j,l}, \quad l \in \{\text{низк}, \text{сред}, \text{выс}\},
		\end{equation}
		где
		\begin{itemize}
			\item \(l\) --- рассматриваемая степень проявления грани;
			\item \(K_l\) --- количество пар <<ситуация + степень>>, относящихся к степени \(l\);
			\item \(M_{j,l}\) --- соответствие разработанному методу для \(j\)-й пары <<ситуация + степень>> при степени проявления \(l\);
			\item \(M_l\) --- среднее совпадение с результатом разработанного метода по всем парам, относящимся к степени проявления \(l\).
		\end{itemize}

		Совпадение с разработанным методом рассчитывалось для всех 90 пар <<ситуация + степень>>, для которых был задан эталонный вариант реакции.
		Средний процент совпадения с разработанным методом составил 45~\%.
		По степеням проявления получены следующие значения: для низкой --- 48~\%, для средней --- 38~\%, для высокой --- 48~\%.

		\includesvgimage
		{research_method_match_heatmap}
		{f}
		{H}
		{0.95\textwidth}
		{Агрегированные данные опроса респондентов}

В рамках этапа 5 выполнен отбор реакций на основании числа голосов
	респондентов.
	Для каждой пары <<ситуация + степень>> отдельно учитывалась каждая
	выбранная реакция: если респондент выбирал две реакции, каждая из них
	получала один голос.
	В качестве результата отбора фиксировалась реакция с наибольшим числом
	голосов.
	Если несколько реакций набирали одинаковое максимальное число голосов,
	они сохранялись как равноправные варианты для дальнейшего использования.

	Всего было рассмотрено 90 пар <<ситуация + степень>>.
	По результатам отбора получено 99 реакций-лидеров: в 81 случае была
	выделена одна реакция, а в 9 случаях две реакции получили одинаковое
	максимальное число голосов.
	В 83 случаях из 90 лидирующая реакция получила не менее 10 голосов из 19,
	то есть была поддержана большинством респондентов.

	Таким образом, для большей части пар <<ситуация + степень>> удалось
	выделить реакцию, которая воспринимается респондентами как характерная
	для заданного значения грани.
	Только 7 пар не получили поддержки большинства.

	Результаты отбора различаются по степеням проявления грани.
	Для низкой степени не менее 10 голосов получила 28 из 30 реакций-лидеров,
	для средней --- 27 из 30, для высокой --- 28 из 30.
	Не менее 14 голосов получили 18 реакций-лидеров для низкой степени, 8
	для средней и 22 для высокой.

	Подробный перечень отобранных реакций приведён в таблице~\ref{tbl:selection_results_by_pair}.
	Номер соответствует номеру ситуации в анкете.

	\begin{landscape}
	\begingroup
	\small
	\begin{longtable}{|p{0.06\linewidth}|p{0.2\linewidth}|p{0.11\linewidth}|p{0.08\linewidth}|p{0.46\linewidth}|}
		\caption{Отобранные реакции по каждой паре <<ситуация + степень>>}
		\label{tbl:selection_results_by_pair} \\
		\hline
		\textbf{Номер} & \textbf{Грань} & \textbf{Степень} & \textbf{Голоса} & \textbf{Отобранные реакции} \\
		\hline
		\endfirsthead
		\caption[]{Отобранные реакции по каждой паре <<ситуация + степень>> (продолжение)} \\
		\hline
		\textbf{Номер} & \textbf{Грань} & \textbf{Степень} & \textbf{Голоса} & \textbf{Отобранные реакции} \\
		\hline
		\endhead
		\hline
		\endfoot
		\endlastfoot
		1 & дружелюбие & низкая & 17/19 & Ты хочешь какую то ерунду! \\
		\hline
		1 & дружелюбие & средняя & 15/19 & Хорошо, помоги \\
		\hline
		1 & дружелюбие & высокая & 19/19 & Я восхищен твоим стремлением помогать людям! \\
		\hline
		2 & общительность & низкая & 16/19 & Не стоит тебе идти на эту вечеринку \\
		\hline
		2 & общительность & средняя & 14/19 & Можно было бы, но сейчас в приоритете другое \\
		\hline
		2 & общительность & высокая & 19/19 & Конечно, я обожаю вечеринки! \\
		\hline
		3 & настойчивость & низкая & 11/19 & Спасибо; Что взял это на себя \\
		\hline
		3 & настойчивость & средняя & 12/19 & Рад за тебя; Что ты готов потратить на это время \\
		\hline
		3 & настойчивость & высокая & 9/19 & Реакции-лидеры (2): Я восхищен твоим усердием!; Я сам хотел сделать ее! \\
		\hline
		4 & уровень активности & низкая & 13/19 & Не стоит тебе идти на эту вечеринку \\
		\hline
		4 & уровень активности & средняя & 12/19 & Рад; Что ты предложил мне это \\
		\hline
		4 & уровень активности & высокая & 19/19 & Конечно, я обожаю вечеринки! \\
		\hline
		5 & поиск острых ощущений & низкая & 12/19 & Реакции-лидеры (2): Кажется, ты не подходишь под ограничения для него; Это не очень хорошая идея кататься на нем \\
		\hline
		5 & поиск острых ощущений & средняя & 9/19 & Рад за тебя; Что ты смог его отыскать! \\
		\hline
		5 & поиск острых ощущений & высокая & 15/19 & Здорово; Что мы сходимся во мнении по поводу лучшего вида аттракционов \\
		\hline
		6 & жизнерадостность & низкая & 13/19 & Реакции-лидеры (2): Всё утекает, как дождевая вода.; Как неприятно, на улице стало неуютно. \\
		\hline
		6 & жизнерадостность & средняя & 10/19 & Мне жаль, что ты застал его \\
		\hline
		6 & жизнерадостность & высокая & 17/19 & Сейчас бы пойти погулять под дождём, внести нотку новизны! \\
		\hline
		7 & доверчивость & низкая & 17/19 & Глупая просьба! Зачем она тебе? \\
		\hline
		7 & доверчивость & средняя & 11/19 & Хорошо; Что она у меня есть \\
		\hline
		7 & доверчивость & высокая & 15/19 & Конечно! Рад; Что спросил! \\
		\hline
		8 & моральность & низкая & 7/19 & Реакции-лидеры (2): Классно; Что ты отдаешь всего себя людям; Это радостная новость \\
		\hline
		8 & моральность & средняя & 12/19 & Реакции-лидеры (2): Жаль; Что ради этого тебе пришлось меньше спать; Это радостная новость \\
		\hline
		8 & моральность & высокая & 16/19 & Я так рад! Теперь ты можешь купить себе все что хочешь \\
		\hline
		9 & альтруизм & низкая & 17/19 & Тебе уже ничего не поможет \\
		\hline
		9 & альтруизм & средняя & 15/19 & Я помогаю людям \\
		\hline
		9 & альтруизм & высокая & 16/19 & Рад; Что ты обратился ко мне за помощью \\
		\hline
		10 & сотрудничество & низкая & 16/19 & Это моя задача и я сам хотел сделать ее! \\
		\hline
		10 & сотрудничество & средняя & 12/19 & Рад за тебя; Что ты готов потратить на это время \\
		\hline
		10 & сотрудничество & высокая & 15/19 & Я восхищен твоей усердностью! \\
		\hline
		11 & скромность & низкая & 16/19 & Да, я всегда был и буду чертовски крут! \\
		\hline
		11 & скромность & средняя & 15/19 & Меня радуют твои слова \\
		\hline
		11 & скромность & высокая & 14/19 & Нет, это ты большой молодец! \\
		\hline
		12 & сочувствие & низкая & 14/19 & Ну ты хотя бы набрался опыта \\
		\hline
		12 & сочувствие & средняя & 12/19 & Я рад; Что ты делишься со мной этим \\
		\hline
		12 & сочувствие & высокая & 18/19 & Мне очень жаль; Что ничего не вышло \\
		\hline
		13 & самоэффективность & низкая & 16/19 & Я не контроллирую ситуацию! \\
		\hline
		13 & самоэффективность & средняя & 10/19 & Мне просто никто не помогает в этом \\
		\hline
		13 & самоэффективность & высокая & 13/19 & Не пытайся мне насолить, я справлюсь \\
		\hline
		14 & аккуратность & низкая & 10/19 & Это место не очень удобное для уборки \\
		\hline
		14 & аккуратность & средняя & 9/19 & Спасибо тебе; Что прибрался у меня! \\
		\hline
		14 & аккуратность & высокая & 16/19 & Да, это стоило всех приложенных мной усилий на уборку \\
		\hline
		15 & исполнительность & низкая & 17/19 & Ничего не выйдет \\
		\hline
		15 & исполнительность & средняя & 9/19 & Вам бы не пришлось просить, если бы вы сами что то для этого сделали \\
		\hline
		15 & исполнительность & высокая & 19/19 & Спасибо за ценную возможность! Я вас не подведу \\
		\hline
		16 & стремление к достижениям & низкая & 16/19 & Вы хотите, чтобы я опозорился там? \\
		\hline
		16 & стремление к достижениям & средняя & 10/19 & На меня будут смотреть люди \\
		\hline
		16 & стремление к достижениям & высокая & 18/19 & Я восхищен вашим предложением! \\
		\hline
		17 & самодисциплина & низкая & 18/19 & У меня не получится \\
		\hline
		17 & самодисциплина & средняя & 17/19 & Другие задачи не дадут мне заняться вашей \\
		\hline
		17 & самодисциплина & высокая & 19/19 & Конечно! Сделаю все сегодня! \\
		\hline
		18 & осторожности & низкая & 16/19 & Что ж, меньше народу больше кислороду \\
		\hline
		18 & осторожности & средняя & 13/19 & Она умерла \\
		\hline
		18 & осторожности & высокая & 18/19 & Мне очень жаль, я понимаю твое горе, держись \\
		\hline
		19 & тревожности & низкая & 13/19 & Я проведу время в хорошей компании \\
		\hline
		19 & тревожности & средняя & 14/19 & Я проведу время в необычной компании \\
		\hline
		19 & тревожности & высокая & 12/19 & О это будет ужасно! \\
		\hline
		20 & гневность & низкая & 14/19 & Спасибо, что сообщил \\
		\hline
		20 & гневность & средняя & 11/19 & Если это шутка, то она невероятно глупая и неуместная! \\
		\hline
		20 & гневность & высокая & 11/19 & Ты только и делаешь что приносишь плохие новости! \\
		\hline
		21 & подавленность & низкая & 10/19 & Реакции-лидеры (2): Здорово, что тебе удалось увидеть его; Сейчас бы пойти погулять под дождём, внести нотку новизны! \\
		\hline
		21 & подавленность & средняя & 16/19 & Мне жаль, что ты застал его \\
		\hline
		21 & подавленность & высокая & 9/19 & Реакции-лидеры (2): Всё утекает, как дождевая вода.; Как неприятно, на улице стало неуютно. \\
		\hline
		22 & чувствительность к оценке & низкая & 12/19 & Реакции-лидеры (2): Прекрасная идея!; Я восхищен вашим предложением! \\
		\hline
		22 & чувствительность к оценке & средняя & 12/19 & Вы ошибаетесь в выборе кандидата для такого мероприятия \\
		\hline
		22 & чувствительность к оценке & высокая & 13/19 & Вы хотите, чтобы я опозорился там? \\
		\hline
		23 & неумеренность & низкая & 12/19 & Хороший торт! \\
		\hline
		23 & неумеренность & средняя & 10/19 & Маловат торт для сюрприза \\
		\hline
		23 & неумеренность & высокая & 13/19 & Нужно было делать торт побольше! \\
		\hline
		24 & уязвимость & низкая & 9/19 & Спасибо, что помог мне осознать это \\
		\hline
		24 & уязвимость & средняя & 10/19 & Мне хотя бы никто не сует палки в колеса \\
		\hline
		24 & уязвимость & высокая & 12/19 & У меня ничего не получается! \\
		\hline
		25 & воображение & низкая & 15/19 & Рад за то, что ты остаешься объективным! \\
		\hline
		25 & воображение & средняя & 12/19 & В некоторых ситуациях фантазия выигрывает у фактов \\
		\hline
		25 & воображение & высокая & 14/19 & В некоторых ситуациях фантазия выигрывает у фактов \\
		\hline
		26 & художенственные интересы & низкая & 16/19 & Ты хочешь купить какую то ерунду! \\
		\hline
		26 & художенственные интересы & средняя & 12/19 & Ты хочешь купить ценную вещь \\
		\hline
		26 & художенственные интересы & высокая & 18/19 & Я восхищен твоим выбором! Она прекрасна! \\
		\hline
		27 & эмоциональность & низкая & 13/19 & Ты говоришь ерунду какую то \\
		\hline
		27 & эмоциональность & средняя & 10/19 & Ты обманываешь \\
		\hline
		27 & эмоциональность & высокая & 18/19 & Я восхищен твоей осознанностью твоих чувств! \\
		\hline
		28 & авантюризм & низкая & 14/19 & Ты собираешься сделать какую то ерунду! \\
		\hline
		28 & авантюризм & средняя & 13/19 & Радует то, что ты достигнешь своего дома в любом случае \\
		\hline
		28 & авантюризм & высокая & 18/19 & Здорово! Ты откроешь для себя новый путь домой! \\
		\hline
		29 & интеллектуальная любознательность & низкая & 16/19 & Ты хочешь заняться какой то ерундой \\
		\hline
		29 & интеллектуальная любознательность & средняя & 16/19 & Я помогу ее тебе решить \\
		\hline
		29 & интеллектуальная любознательность & высокая & 19/19 & Я помогу ее тебе решить \\
		\hline
		30 & либерализм & низкая & 15/19 & Ты хочешь заниматься какой то ерундой! \\
		\hline
		30 & либерализм & средняя & 10/19 & Рад за то, что ты нашел свое место! \\
		\hline
		30 & либерализм & высокая & 14/19 & Реакции-лидеры (2): Вау! Ты собираешься выступить против такого сильного врага!; Восхищает твое желание работать для важной цели! \\
		\hline
	\end{longtable}
	\endgroup
	\end{landscape}

	Полученные результаты отбора подтверждают, что наиболее определённо
	респонденты распознают реакции, соответствующие высокой степени проявления
	грани.
	Средняя степень чаще воспринимается неоднозначно.
	Реакции, получившие единичные голоса, не должны включаться в основной
	набор правил.

В рамках этапа 6 сформулированы рекомендации по модификации набора
	правил на основании интерпретации мнений респондентов:
\begin{itemize}
	\item выбрать 83 сценария, поддержанных не менее, чем 10 респондентами, и для каждой пары <<ситуация + степень>> из этих 83 случаев провести корректировку правил;
	\item правила, результат применения которых в виде реакций не получил поддержку большинства (не менее 10 респондентов) следует переработать, и затем обновленные реакции повторно подвергнуть исследовательской процедуре;
	\item правила, результат применения которых в виде реакций не получил поддержку даже 2--3 респондентов, следует дополнить и переработать, и затем обновленные реакции повторно подвергнуть исследовательской процедуре.
\end{itemize}


\phantomsection
\section*{Выводы}
\addcontentsline{toc}{section}{Выводы}

В исследовательском разделе описана анкета для оценки различимости реакций
робота по 30 граням Большой пятёрки и трём степеням проявления каждой грани.
По результатам опроса средняя согласованность составила 50~\%, а в 12 из 90 случаев согласованность была не ниже 75~\%.
Средний процент совпадения с разработанным методом составило 45~\%.
Анализ отдельных реакций показал, что 33~\% реакций были выбраны не менее чем 10 респондентами, поэтому часть реакций воспринимается устойчиво даже при расхождении полных наборов ответов.
Это показывает, что предложенный набор ситуаций и реакций позволяет различать проявления части граней, но значительная часть граней требует дальнейшего уточнения.
Также описана исследовательская процедура, согласно которой выполнен отбор реакций, обусловленных правилами: 99 реакций, основанных на сценариях реагирования робота в коммуникации, рекомендуются к применению для заданых степеней проявления соответствующих граней.
Предложенные метод и исследовательская процедура рекомендуются к применению для роботов-компаньонов.